% Layout and general macros, taken from the owner's LaTeX master (revtex4-2; colors unchanged).
% Notation macros of this project go in notation.tex, not here.

% cancel automatic upper case in the section title
\makeatletter
\def\@hangfrom@section#1#2#3{\normalsize\@hangfrom{#1#2}#3}%\MakeTextUppercase{#3}}%
\def\@hangfroms@section#1#2{\normalsize#1#2}%\MakeTextUppercase{#2}}%
\makeatother

\RequirePackage[usenames, dvipsnames, x11names]{xcolor}
% Useful packages
\usepackage{amsmath,amsfonts,amssymb,bm}
\usepackage{graphicx,color,soul}
\usepackage[normalem]{ulem}
\usepackage{dsfont}
\usepackage{enumerate}
\usepackage{multirow}
\usepackage{longtable}
\usepackage[colorlinks=true, allcolors=blue]{hyperref}

\hypersetup{
    colorlinks=true,
    citecolor=Green,
    linkcolor=Red,
    urlcolor=NavyBlue
}
\graphicspath{{figs/}}

% new command = nc
\newcommand{\nc}{\newcommand}
% dirac notation
\nc{\ket}[1]{|#1\rangle}
\nc{\bra}[1]{\langle#1|}
\nc{\ketbra}[2]{|#1\rangle\!\langle#2|}
\nc{\braket}[2]{\langle#1|#2\rangle}
\nc{\braoprket}[3]{\langle#1|#2|#3\rangle}
\nc{\opr}[1]{\operatorname{#1}}
\nc{\avg}[1]{\langle#1\rangle}
\nc{\ketbrasame}[1]{|#1\rangle\!\langle#1|}
\nc{\tr}{\opr{tr}}
\nc{\E}{\mathbb{E}}
\nc{\var}{\opr{Var}}
\nc{\up}{\uparrow}
\nc{\dn}{\downarrow}
\nc{\bd}[1]{\boldsymbol{#1}}
\nc{\zjx}[1]{\textcolor{violet}{\textbf{#1}}}
\nc{\jx}[1]{\textcolor{violet}{\textbf{[JX: #1]}}}
\nc{\red}[1]{\textcolor{red}{#1}}

\newcommand{\diag}[2]{\vcenter{\hbox{\includegraphics[height=#2]{#1}}}}
\newcommand{\figref}[1]{Fig.\,\ref{#1}}
\newcommand{\figsref}[1]{Figs.\,\ref{#1}}
\newcommand{\secref}[1]{Sec.\,\ref{#1}}
\newcommand{\eqnref}[1]{Eq.\,\eqref{#1}}
\newcommand{\eqnsref}[1]{Eqs.\,\eqref{#1}}
\newcommand{\appref}[1]{Appendix~\ref{#1}}
\renewcommand{\Im}{\mathop{\mathrm{Im}}}
\renewcommand{\Re}{\mathop{\mathrm{Re}}}
\nc{\sgn}{\operatorname{sgn}}
\nc{\HC}{\mathrm{H.c.}}
\nc{\Tr}{\mathrm{Tr}}
\nc{\ba}{\mathbf{a}}
\nc{\bb}{\mathbf{b}}
\nc{\bk}{\mathbf{k}}
\nc{\bq}{\mathbf{q}}
\nc{\br}{\mathbf{r}}
\nc{\bQ}{\mathbf{Q}}

% ---------------------------------------------------------------------------
% Review marks (colors as in the master). In each batch the owner's own material is black,
% details filled in by the agent are orange, and the agent's questions are blue. After approval,
% rd mynotes accept strips the orange marks; it refuses while any blue question is left.
%   \orange{...}                               detail filled in by the agent (sentence, short formula)
%   \begin{orangeblock} ... \end{orangeblock}  added passage, may contain displayed equations;
%                                              it does not start a new paragraph by itself
%   \cmt{...}                                  question or doubt from the agent
%   \jx{...}                                   the owner's own remark; accept leaves it alone
%   \tbc                                       "(to be completed)" on a line of its own
\nc{\cmt}[1]{\textcolor{blue}{\textbf{[#1]}}}
\DeclareRobustCommand{\orange}[1]{\textcolor{BurntOrange}{\textbf{\boldmath #1}}}
\newenvironment{orangeblock}{\color{BurntOrange}\bfseries\boldmath\ignorespaces}{\ignorespacesafterend}
\nc{\tbc}{\par\noindent{\color{black}\textbf{(to be completed)}}\par}

% ---------------------------------------------------------------------------
% Notation table: each \notation{symbol}{meaning}{defined in} in notation.tex is one row;
% \printnotation at the end prints the table (nothing while there are no entries).
\makeatletter
\newif\if@rdnotation
\newcommand\rd@notationrows{}
\newcommand{\notation}[3]{\global\@rdnotationtrue\g@addto@macro\rd@notationrows{#1 & #2 & #3\\}}
\newcommand{\printnotation}{%
  \if@rdnotation
    \section*{Notation}
    \begin{longtable}{l p{0.58\textwidth} l}
      \hline
      Symbol & Meaning & Defined in\\
      \hline
      \rd@notationrows
      \hline
    \end{longtable}
  \fi}
\makeatother
